\section{维数正规化下非局域算符的重正化}

现在我们来考虑出现在部分子物理 LaMET 方法中的非局域算符（如式~(\ref{Eq:nonlocalop})）的重正化。在有效理论中，把非局域威尔逊线替换为两个辅助“重夸克”场之积后，非局域算符成为定域复合算符的乘积，
\begin{equation}
 O(z_2, z_1) = \overline{j}(z_2) j(z_1),
\end{equation}
该算符可以通过
 \begin{equation}
        O(z_2,z_1)= Z_{\bar j}Z_{j} O_R(z_2,z_1)
\end{equation}
进行乘法重正化，且在微扰论所有阶成立，其中 $Z_{j}=Z_q^{1/2} Z_Q^{1/2} Z_{V_j}$，而 $Z_q, Z_Q, Z_{V_j}$ 分别是轻夸克、重夸克和顶点的重正化常数。

于是重正化的算符成为 $O_R=Z_{\bar j}^{-1}Z_{j}^{-1}O(z_2,z_1)$。在对“重夸克”场积分之后，我们得到
\begin{equation}\label{Eq:Orenorm}
        O_R = Z_{\bar j}^{-1}Z_{j}^{-1} \overline{\psi}(z_2) \Gamma L(z_2,z_1)\psi(z_1),
\end{equation}
其中右边的场都是裸场。文献~\cite{Ji:2015jwa} 已经证明，准 PDF 算符只需要虚图的重正化，这与式~(\ref{Eq:Orenorm}) 的结果等价：对算符 $O$ 我们只需要轻重夸克流的重正化。特别是，准 PDF 的单圈重正化因子
\beq\label{Eq:1loopZfac}
Z_O=1+\frac{\alpha_s }{\pi\epsilon}
\eeq
与式~(\ref{eq:hl}) 中轻重夸克流的反常量纲相一致。两圈阶的情形也是如此~\cite{Ji:2015jwa}。
